\cotefile{15}
\scanpage{27}
\mspage{\handwritten{48}}
% archive sheet number: 26

L'élément $\handwritten{\int_\gamma}(\handwritten{\sum_k f}_k\,dx_k)$ de E ne dépend donc que de $\handwritten{a}$ et $\handwritten{z}$ et non de la route $\handwritten{\gamma}$ ; $\handwritten{a}$ restant fixe, posons $\handwritten{g}(\handwritten{z})=\handwritten{\int_\gamma}(\handwritten{\sum_k f}_k\,dx_k)$. Or, pour tout $\handwritten{z_0}\handwritten{\in}A$, il existe un polycylindre P de centre $\handwritten{z_0}$, contenu dans A, et une fonction holomorphe $\handwritten{u}$ dans P telle que $\handwritten{u}(\handwritten{z_0})=\handwritten{g}(\handwritten{z_0})$ et $d\handwritten{u}=\handwritten{\sum_k f}_k\,dx_k$ dans P ; si pour $\handwritten{z}\handwritten{\in}P$, $\handwritten{\sigma}$ est la route linéaire $t\handwritten{\to z_0}+t(\handwritten{z}-\handwritten{z_0})$ $(t\handwritten{\in[}0,1\handwritten{]})$, on a donc $\handwritten{u}(\handwritten{z})-\handwritten{u}(\handwritten{z_0})=\handwritten{\int_\sigma}(\handwritten{\sum_k f}_k\,dx_k)=\handwritten{g}(\handwritten{z})-\handwritten{g}(\handwritten{z_0})$, d'où $\handwritten{g}(\handwritten{z})=\handwritten{u}(\handwritten{z})$ dans P, ce qui démontre la proposition.

En particulier, si $n=1$, on voit que toute fonction holomorphe $\handwritten{f}$ dans un domaine \uline{contractile} $A\handwritten{\subset}K$ admet une primitive $\handwritten{g}$ holomorphe dans A, c'est-à-dire une fonction dont la dérivée $\handwritten{g}'=\handwritten{f}$ dans A ; en outre, si $\handwritten{g}_1$ et $\handwritten{g}_2$ sont deux primitives de $\handwritten{f}$ dans A, leur différence est \deleted{\unsure{même}} constante dans un voisinage d'un point quelconque de A, donc constante dans A puisque A est connexe.

Une variante de la démonstration du th.1 et de la prop.2 consiste à établir d'abord un principe général de prolongement analytique, de la façon suivante :

\uline{Etant donné un domaine} $A\handwritten{\subset}K^n$, \uline{on suppose donnée en chaque point} $\handwritten{z}\handwritten{\in}A$, \uline{une famille} $\handwritten{\Phi}(\handwritten{z})$ \uline{de polycylindres ouverts contenus dans A et de centre} $\handwritten{z}$, \uline{et pour tout} $P\handwritten{\in\Phi}(\handwritten{z})$, \uline{une fonction holomorphe} $\handwritten{u}_{\unsure{P|p}}$ \uline{définie dans} P ; \uline{on suppose que pour tout} $\handwritten{x\in}P\handwritten{\in\Phi}(\handwritten{z})$ \uline{il existe} $Q\handwritten{\in\Phi}(\handwritten{x})$ \uline{tel que} $\handwritten{u}_Q(\handwritten{x})=\handwritten{u}_P(\handwritten{x})$ ; \uline{on suppose en outre que si} $P\handwritten{\in\Phi}(\handwritten{z})$, $Q\handwritten{\in\Phi}(\handwritten{z}')$ \uline{ont une intersection non vide et si} $\handwritten{u}_P(\handwritten{x})=\handwritten{u}_Q(\handwritten{x})$ \uline{en un point de} $P\handwritten{\cap}Q$, \uline{alors} $\handwritten{u}_P$ \uline{et} $\handwritten{u}_Q$ \uline{sont identiques dans} $P\handwritten{\cap}Q$ ; \uline{enfin}, \uline{si}
\scanpage{29}
\mspage{47}
% archive sheet number: 28

en particulier $\handwritten{\int_\gamma}d\handwritten{g}=0$ le long de tout \uline{circuit} dans A.

PROPOSITION 2.- \uline{Soit} $A\handwritten{\subset}K^n$ \uline{un domaine} contractile, \uline{et soit} $\handwritten{\sum_{k=1}^n f}_k\,dx_k$ \uline{une expression différentielle fermée à coefficients holomorphes dans} A. \uline{Il existe alors une fonction} $\handwritten{g}$ \uline{holomorphe dans} A \uline{et telle que} $d\handwritten{g}=\handwritten{\sum_{k=1}^n f}_k\,dx_k$.

Nous utiliserons le lemme suivant :

\uline{Lemme 1}.- \uline{Quels que soient les points} $\handwritten{a}$, $\handwritten{b}$ \uline{d'un domaine} $A\handwritten{\subset}K^n$, \uline{il existe une route a-chevée linéaire par morceaux, contenue dans} A \uline{et ayant pour extrémités} $\handwritten{a}$ \uline{et} $\handwritten{b}$.

Il suffit de prouver que l'ensemble B des extrémités des routes linéaires par morceaux contenues dans A et d'origine A, est à la fois ouvert et fermé. Or, si $\handwritten{c\in}\overline{B}$, il existe une boule S de centre $\handwritten{c}$ contenue dans A et contenant une extrémité $\handwritten{b}$ d'une route $\handwritten{\gamma}$ d'origine $\handwritten{a}$ contenue dans A. Comme le segment d'extrémités $\handwritten{b}$ et $\handwritten{c}$ est contenu dans S, donc dans A, $\handwritten{c}$ est l'extrémité d'une route linéaire par morceaux, juxta-position de $\handwritten{\gamma}$ et d'une route linéaire. De même, tout point $\handwritten{x\in}S$ est extrémité d'une route linéaire par morceaux, juxtaposition de $\handwritten{\gamma}$ et d'une route linéaire, donc $S\handwritten{\subset}C$, ce qui prouve que C est ouvert et fermé. Ce lemme étant démontré, soit $\handwritten{a}$ un point quelconque de A, $\handwritten{z}$ un point quelconque de A, $\handwritten{\gamma}$ et $\handwritten{\gamma}'$ deux routes contenues dans A d'origine $\handwritten{a}$ et d'extrémité $\handwritten{z}$ (il en existe en vertu du lemme). On a $\handwritten{\int_\gamma}(\handwritten{\sum_k f}_k\,dx_k)=\handwritten{\int_{\gamma'}}(\handwritten{\sum_k f}_k\,dx_k)$. En effet, si $\handwritten{\gamma}$ est définie dans un intervalle $\handwritten{[}b,c\handwritten{]}$, on peut, en remplaçant $\handwritten{\gamma}'$ par une route équivalente, supposer que $\handwritten{\gamma}'$ est définie dans un intervalle de la forme $\handwritten{[}c,d\handwritten{]}$ ; on a $\handwritten{\gamma}(c)=\handwritten{\gamma}'(d)=\handwritten{z}$ ; on peut donc définir $\handwritten{\gamma}''$ dans $\handwritten{[}b,d\handwritten{]}$, en posant $\handwritten{\gamma}''(t)=\handwritten{\gamma}(t)$ dans $\handwritten{[}b,c\handwritten{]}$ et $\handwritten{\gamma}''(t)=\handwritten{\gamma}'(c+d-t)$ dans $\handwritten{[}c,d\handwritten{]}$ ; comme $\handwritten{\gamma}''(d)=\handwritten{\gamma}'(c)=\inserted{\handwritten{\unsure{a|z}=}}\handwritten{\gamma}''(b)$, $\handwritten{\gamma}''$ est un \uline{circuit}, et notre assertion résulte du th.1.
\scanpage{31}
\mspage{46}
% archive sheet number: 30

COROLLAIRE.- \uline{Soient} $\handwritten{\gamma}$ \uline{et} $\handwritten{\gamma}'$ \uline{deux routes achevées à valeurs dans} A, \uline{ayant même origine} $\handwritten{a}$ \uline{et même extrémité} $\handwritten{b}$, \uline{et telles qu'il existe une déformation de} $\handwritten{\gamma}$ \uline{en} $\handwritten{\gamma}'$ \uline{dans} A \uline{qui laisse fixes} $\handwritten{a}$ \uline{et} $\handwritten{b}$. \uline{Alors, si l'expression différentielle} $\handwritten{\sum_k f}_k\,dx_k$ \uline{satisfait aux conditions} (1), \uline{on a} $\handwritten{\int_\gamma}(\handwritten{\sum_k f}_k\,dx_k)=\handwritten{\int_{\gamma'}}(\handwritten{\sum_k f}_k\,dx_k)$.

On peut supposer que $\handwritten{\gamma}$ et $\handwritten{\gamma}'$ sont définies dans le même intervalle $I=\handwritten{[\alpha,\beta]}$ ; soit $\handwritten{\gamma}^{\unsure{o|0|\circ}}$ l'opposé de $\handwritten{\gamma}$, et $\handwritten{\gamma}_1(t)=\handwritten{\gamma}_o(t-\handwritten{\beta}+\handwritten{\alpha})$ un chemin équivalent à $\handwritten{\gamma}^o$ défini dans $I+(\handwritten{\beta}-\handwritten{\alpha})$. Les routes $\handwritten{\gamma\vee\gamma}_1$ et $\handwritten{\gamma}'\handwritten{\vee\gamma}_1$ sont des circuits dans A et peuvent évidemment être déformés l'un en l'autre dans A en vertu des hypothèses ; en leur appliquant le th.1, on a aussitôt le corollaire.

\unsure{DÉFINITION|DÈFINITION} 5.- \uline{On dit qu'un domaine} $A\handwritten{\subset}K^n$ \uline{est contractile si tout lacet dans} A \uline{peut être déformé dans} A \uline{en un point de} A.

Par exemple, tout domaine \uline{étoilé} par rapport à un point est contractile

Une expression différentielle $\handwritten{\sum_{k=1}^n f}_k\,dx_k$ à coefficients holomorphes dans un domaine A est dite \uline{fermée} si elle satisfait aux conditions (1). On voit donc que :

\unsure{THÈORÈME|THÉORÈME|THÈORÉME|THÉORÉME} 2 (Cauchy).- \uline{Dans un domaine contractile l'intégrale d'une expression différentielle fermée} (à coefficients holomorphes) \uline{le long d'un circuit quelconque est nulle}.

On peut évidemment, pour le cas des variables \uline{réelles}, généraliser le th.1 au cas où les $\handwritten{f}_k$ sont seulement supposées continument différentiables ; est-ce que cela en vaut la peine ?

\handwritten{\unsure{3|2}}. \uline{Primitives d'une fonction holomorphe}.

Soit $A\handwritten{\subset}K^n$ un domaine \uline{quelconque}, $\handwritten{g}$ une fonction holomorphe dans A ; si on pose $d\handwritten{g}=\handwritten{\sum_{k=1}^n f}_k\,dx_k$, il est immédiat que cette expression différentielle est \uline{fermée} ; en outre, en vertu des définitions, le long d'une route $\handwritten{\gamma}$ définie dans $I=\qquad$, on a $\handwritten{\int_\gamma}d\handwritten{g}=\handwritten{\int_\gamma}(\handwritten{\sum_k f}_k\,dx_k)=\handwritten{g}(\handwritten{\gamma}(\handwritten{\beta}))-\handwritten{g}(\handwritten{\gamma}(\handwritten{\alpha}$
% The source leaves the space after I= blank. The right edge ends after alpha; no closing parentheses have been supplied.
\scanpage{33}
\mspage{44}
% archive sheet number: 32

chacune des fonctions $\handwritten{f}_k$ soit égale à une série entière sommable par rapport aux $x_j-a_{ij}$. Il existe un nombre $\handwritten{\rho}>0$ tel que, pour tout $\handwritten{x\in\unsure{\varphi|\phi}}(I\handwritten{\times}J)$, le polycylindre ouvert de centre $\handwritten{x}$ et de rayons $\handwritten{\rho}$ soit contenu dans un $P_i$ au moins ; car si $\handwritten{\delta}(\handwritten{x})$ désigne la plus grande distance de $\handwritten{x}$ aux frontières de ceux des $P_i$ qui contiennent $\handwritten{x}$, on vérifie aussitôt que $\handwritten{\delta}(\handwritten{x})$ est continue dans $\handwritten{\varphi}(I\handwritten{\times}J)$ et est $>0$ en tout point, donc atteint son minimum $\handwritten{\rho}$ en un point de $\handwritten{\varphi}(I\handwritten{\times}J)$, ce qui implique $\handwritten{\delta}(\handwritten{x})\handwritten{\geq\rho}>0$ pour tout $\handwritten{x\in\varphi}(I\handwritten{\times}J)$. Par suite, pour tout $\handwritten{x\in\varphi}(I\handwritten{\times}J)$, chacune des fonctions $\handwritten{f}_k$ est égale à une série entière sommable (par rapport aux $z_j-x_j$) dans le polycylindre ouvert de centre $\handwritten{x}$ et de rayon $\handwritten{\rho}$ (chap.I, \handwritten{\S}1, prop.\handwritten{\unsure{4}}). En vertu de la continuité uniforme de $\handwritten{\varphi}$, il existe $\handwritten{\varepsilon}>0$ tel que les relations $|t-t'|\handwritten{\leq\varepsilon}$, $|s-s'|\handwritten{\leq\varepsilon}$ entraînent $\handwritten{\|\varphi}(t,s)-\handwritten{\varphi}(t',s')\handwritten{\|\leq\frac{\rho}{4}}$. Soit $(t_{\handwritten{\alpha}})$ $(0\handwritten{\leq\alpha\leq}p)$ une suite croissante de points de I tels que $t_0=a$, $t_p=b$ et $t_{\handwritten{\alpha}+1}-t_{\handwritten{\alpha}}\handwritten{\leq\varepsilon}$ pour $0\handwritten{\leq\alpha\leq}p-1$ ; soit de même $(s_{\handwritten{\beta}})$ $(0\handwritten{\leq\beta\leq}q)$ une suite croissante de points de J tels que $s_0=c$, $s_q=d$ et $s_{\handwritten{\beta}+1}-s_{\handwritten{\beta}}\handwritten{\leq\varepsilon}$ pour $0\handwritten{\leq\beta\leq}q-1$. Désignons par $\handwritten{\gamma_\beta}$ le circuit défini par
\[
\begin{aligned}
\handwritten{\gamma_\beta}(t)&=\handwritten{\varphi}(t_{\handwritten{\alpha}},s_{\handwritten{\beta}})+\frac{t-t_{\handwritten{\alpha}}}{t_{\handwritten{\alpha}+1}-t}\bigl(\handwritten{\varphi}(t_{\handwritten{\alpha}+1},s_{\handwritten{\beta}})-\handwritten{\varphi}(t_{\handwritten{\alpha}},s_{\handwritten{\beta}})\bigr)\quad\text{pour }t_{\handwritten{\alpha}}\handwritten{\leq}t\handwritten{\leq}t_{\handwritten{\alpha}+1}\\
&0\handwritten{\leq\alpha\leq}p-1,\quad1\handwritten{\leq\beta\leq}q-1
\end{aligned}
\]
; on posera en outre $\handwritten{\gamma}_0=\handwritten{\gamma}$, $\handwritten{\gamma}_q=\handwritten{\gamma}'$. Tout revient à prouver que l'on a, pour $0\handwritten{\leq\beta\leq}q-1$, $\handwritten{\int_{\gamma_\beta}}(\handwritten{\sum_k f}_k\,dx_k)=$ $=\handwritten{\int_{\gamma_{\beta+1}}}(\handwritten{\sum_k f}_k\,dx_k)$. Or, pour chaque indice $\handwritten{\alpha}$ tel que $0\handwritten{\leq\alpha\leq}p-1$, les points $\handwritten{\gamma_\beta}(t)$ et $\handwritten{\gamma_{\beta+1}}(t)$ pour $t_{\handwritten{\alpha}}\handwritten{\leq}t\handwritten{\leq}t_{\handwritten{\alpha}+1}$ appartiennent tous à un même polycylindre ouvert $Q_{\handwritten{\alpha},\handwritten{\beta}}$ de centre un point de $\handwritten{\varphi}(I\handwritten{\times}J)$ et de rayons $\handwritten{\rho}$ (fig.1) ; on sait d'autre part qu'il existe dans $Q_{\handwritten{\alpha\beta}}$
\scanpage{35}
\mspage{42}
% archive sheet number: 34

convenable $s\handwritten{\to}s+a$ et la déformation $(t,s)\handwritten{\to\unsure{\varphi|\phi}}_2(t,s+a)$ de $\handwritten{\gamma}_1$ en $\handwritten{\gamma}_2$ ; alors si $\handwritten{\varphi}$ est la fonction égale à $\handwritten{\varphi}_1$ dans \handwritten{\unsure{\textasciigrave|\textbackslash}}$I\handwritten{\times}J_1$ et à $\handwritten{\varphi}_2$ dans $I\handwritten{\times}J_2$ (ce qui a un sens, puisque les deux fonctions coïncident sur $I\handwritten{\times}(J_1\handwritten{\cap}J_2)$), il est clair que $\handwritten{\varphi}$ est continue dans $I\handwritten{\times}(J_1\handwritten{\cup}J_2)$ et est une déformation de $\handwritten{\gamma}_0$ en $\handwritten{\gamma}_2$ dans A. Comme en outre tout chemin est déformable en lui-même (en prenant $\handwritten{\varphi}(t,s)=\handwritten{\gamma}(t)$ pour tout $s\handwritten{\in}J$), la relation ``$\handwritten{\gamma}$ est déformable en $\handwritten{\gamma}'$ dans A'' est une \uline{relation d'équivalence} entre les chemins définis dans I et à valeurs dans A.

\uline{Exemple}.- On dit qu'un domaine $A\handwritten{\subset\mathbb{R}}^n$ est \uline{étoilé} par rapport à un point $\handwritten{a\in}A$ si pour tout $\handwritten{x\in}A$, le segment fermé d'extrémités $\handwritten{a}$ et $\handwritten{x}$ est contenu dans A. Dans un domaine étoilé par rapport à A, \uline{tout chemin} (resp. \uline{tout lacet}) $\handwritten{\gamma}$ \uline{est déformable en le point} $\handwritten{a}$ par la déformation $\handwritten{\varphi}(t,s)=\handwritten{a}+(1-s)(\handwritten{\gamma}(t)-\handwritten{a})$ (pour $0\handwritten{\leq}s\handwritten{\leq}1$). Deux chemins (resp. lacets) quelconques sont donc déformables l'un en l'autre dans A.

\uline{Remarque}.- Si $\handwritten{\gamma}$ est une juxta-position de deux chemins $\handwritten{\gamma}_1$, $\handwritten{\gamma}_2$, et $\handwritten{\varphi}$ une déformation de $\handwritten{\gamma}$ en $\handwritten{\gamma}'$ dans A, $\handwritten{\gamma}'$ est juxtaposition des deux chemins $\handwritten{\gamma}'_1$, $\handwritten{\gamma}'_2$ déformés de $\handwritten{\gamma}_1$, $\handwritten{\gamma}_2$ par $\handwritten{\varphi}$. Mais on aura soin d'observer que si $\handwritten{\gamma}_1$ et $\handwritten{\gamma}_2$ peuvent être \uline{chacun} déformés \margin{\handwritten{\unsure{2|Z}}}dans A en deux nouveaux chemins $\handwritten{\gamma}'_1$, $\handwritten{\gamma}'_2$, tels que l'extrémité de $\handwritten{\gamma}'_1$ soit l'origine de $\handwritten{\gamma}'_2$, il ne s'ensuit nullement que $\handwritten{\gamma}$ puisse être déformé dans A en la juxtaposition de $\handwritten{\gamma}'_1$ et de $\handwritten{\gamma}'_2$.

PROPOSITION 1.- \uline{Tout chemin achevé dans} A \uline{peut être déformé dans} A \uline{en une route linéaire. par morceaux ayant mêmes extrémités}.

En effet, soit $\handwritten{\gamma}$ un chemin a chevé à valeurs dans A, défini dans l'intervalle compact $I=\handwritten{[}a,b\handwritten{]}$. L'ensemble $\handwritten{\gamma}(I)$ est compact et contenu dans A ; la distance $\handwritten{\rho}$ de cet ensemble à $\handwritten{\complement}A$ est donc $\handwritten{>}0$.
\scanpage{37}
\mspage{\handwritten{31}}
% archive sheet number: 36

\begin{center}
CHAPITRE II (Etat 1)

FONCTIONS HOLOMORPHES D'UNE VARIABLE COMPLEXE.

\handwritten{\S}1. Le principe du prolongement analytique.
\end{center}

\uline{1}. \uline{Prolongement analytique .d'une fonction holomorphe}.

Dans ce \handwritten{\S} ; K désignera \uline{le corps des nombres réels ou le corps des nombres complexes}, E un espace normé complet sur K.

PROPOSITION 1.- \uline{Dans l'espace} $K^p$, \uline{soient} P \uline{et} Q \uline{deux polycylindres ouverts, de centres} $\handwritten{a}$ \uline{et} $\handwritten{b}$, \uline{tels que} $P\handwritten{\cap}Q$ \uline{ne soit pas vide}. \uline{Soient} $\handwritten{f}(\handwritten{x})=\handwritten{\sum c}_{n_1n_2..n_p}(x_1-a_1)^{n_1}\ldots(x_p-a_p)^{n_p}$ \uline{la somme d'une série entière par rapport aux} $x_i-a_i$, \uline{sommable dans} P, \uline{et} $\handwritten{g}(\deleted{\handwritten{\unsure{g|y|z}}}\inserted{\handwritten{x}})=$ $=\handwritten{\sum d}_{n_1n_2..n_p}(x_1-b_1)^{n_1}\ldots(x_p-b_p)^{n_p}$ \uline{la somme d'une série entière par rapport aux} $x_i-b_i$, \uline{sommable dans}\unsure{-} Q. \uline{S'il existe un ensemble ouvert non vide} $U\handwritten{\subset}P\handwritten{\cap}Q$ \uline{tel que} $\handwritten{f}(\handwritten{x})=\handwritten{g}(\handwritten{x})$ \uline{pour tout} $\handwritten{x\in}U$, \uline{on a} $\handwritten{f}(\handwritten{x})=\handwritten{g}(\handwritten{x})$ \uline{pour} tout $\handwritten{x\in}P\handwritten{\cap}Q$.

Soit en effet $\handwritten{x}$ un point de U et $\handwritten{y}$ un point quelconque de $P\handwritten{\cap}Q$ ; le segment fermé d'extrémités $\handwritten{x}$ et $\handwritten{y}$ appartient tout entier à $P\handwritten{\cap}Q$, qui est convexe ; la fonction $\handwritten{h}(t)=\handwritten{f}(\handwritten{x}+t(\handwritten{y}-\handwritten{x}))-\handwritten{g}(\handwritten{x}+t(\handwritten{y}-\handwritten{x}))$ est donc une fonction holomorphe de la variable réelle t dans un intervalle ouvert I contenant $\handwritten{[}0,1\handwritten{]}$. Désignons par $\handwritten{\rho}$ la borne supérieure de l'ensemble A des t tels que $0\handwritten{\leq}t\handwritten{\leq}1$ et que $\handwritten{h}(s)=0$ pour $0\handwritten{\leq}s\handwritten{\leq}t$ ; on a par hypothèse $\handwritten{\rho}>0$, et il est clair par continuité que $\handwritten{\rho\in}A$ ; nous allons montrer que $\handwritten{\rho}=1$ ce qui démontrera la proposition. Dans le cas contraire, la fonction $\handwritten{h}(t)$ serait égale à la somme d'une série entière en $t-\handwritten{\rho}$, convergente dans un intervalle ouvert de centre $\handwritten{\rho}$ ; mais on a $\handwritten{h}(t)=0$ pour tout t tel que $0\handwritten{\leq}t\handwritten{\leq\rho}$, donc d'après le principe des zéros isolés, les coefficients de cette série en $t-\handwritten{\rho}$ sont tous nuls on a donc aussi $\handwritten{h}(t)=0$ pour $\handwritten{\rho}<t<\handwritten{\rho}+\handwritten{\alpha}$, où $\handwritten{\alpha}>0$ ; mais cela est
\scanpage{39}
\mspage{29}
% archive sheet number: 38

\[
\begin{aligned}
\frac{\handwritten{\partial}}{\handwritten{\partial}x_k}\left(\frac{\handwritten{\partial u}_j}{\handwritten{\partial}x_i}\right)&=\frac{\handwritten{\partial}U_{ij}}{\handwritten{\partial}x_k}(\handwritten{x};\handwritten{u}(\handwritten{x}))+\handwritten{\sum_{h=1}^q}\frac{\handwritten{\partial}U_{ij}}{\handwritten{\partial}y_h}(\handwritten{x};\handwritten{u}(\handwritten{x}))\mathbin{.}\frac{\handwritten{\partial u}_h}{\handwritten{\partial}x_k}
&& (1\handwritten{\leq}i\handwritten{\leq}p,\\
&=\frac{\handwritten{\partial}U_{ij}}{\handwritten{\partial}x_k}(\handwritten{x};\handwritten{u}(\handwritten{x}))+\handwritten{\sum_{h=1}^q}\frac{\handwritten{\partial}U_{ij}}{\handwritten{\partial}y_h}(\handwritten{x};\handwritten{u}(\handwritten{x}))U_{kh}(\handwritten{x};\handwritten{u}(\handwritten{x})).
&&1\handwritten{\leq}k\handwritten{\leq}p,\\
&&&1\handwritten{\leq}j\handwritten{\leq}q)
\end{aligned}
\]

Ecriva\handwritten{-}nt que les expressions ainsi obtenues pour $\frac{\handwritten{\partial}}{\handwritten{\partial}x_k}(\frac{\handwritten{\partial u}_j}{\handwritten{\partial}x_i})$ et $\frac{\handwritten{\partial}}{\handwritten{\partial}x_i}(\frac{\handwritten{\partial u}_j}{\handwritten{\partial}x_k})$ sont égales au point $\handwritten{x}$, puis faisant $\handwritten{x}=\handwritten{x}_0$ et tenant compte de (23) il vient
\[
\begin{aligned}
&\frac{\handwritten{\partial}U_{ij}}{\handwritten{\partial}x_k}(\handwritten{x}_0;\handwritten{y}_0)+\handwritten{\sum_{h=1}^q}\frac{\handwritten{\partial}U_{ij}}{\handwritten{\partial}y_h}(\handwritten{x}_0;\handwritten{y}_0)U_{kh}(\handwritten{x}_0;\handwritten{y}_0)=\\
&=\frac{\handwritten{\partial}U_{kj}}{\handwritten{\partial}x_i}(\handwritten{x}_0;\handwritten{y}_0)+\handwritten{\sum_{h=1}^q}\frac{\handwritten{\partial}U_{kj}}{\handwritten{\partial}y_h}(\handwritten{x}_0;\handwritten{y}_0)U_{jh}(\handwritten{x}_0;\handwritten{y}_0)
\end{aligned}
\]
Mais comme par hypothèse, $(\handwritten{x}_0,\handwritten{y}_0)$ est un point arbitraire de A, on obtient bien (24).

$2^\circ$ Les conditions sont \uline{suffisantes}. Considérons l'équation différentielle
\[
\handwritten{y}'=\underline{U}(\handwritten{x}_0+t(\handwritten{x}-\handwritten{x}_0);\handwritten{y})\mathbin{.}(\handwritten{x}-\handwritten{x}_0)\tag*{(25)}
\]

où $\underline{U}$ est la matrice $(U_{ij})$, $\handwritten{x}_0$ et $\handwritten{x}$ sont fixés dans B, et $\handwritten{y}'=d\handwritten{y}/dt$, t variant dans un voisina ge de 0 dans K. Il résulte de la formation de la série majorante dans la démonstration du th.2 que, dès que $\handwritten{x}-\handwritten{x}_0$ est assez petit, l'équation (25) admet une solution t une seule $\handwritten{v}(t;\handwritten{x}_0,\handwritten{x},\handwritten{y}_0)$ holomorphe pour $|t|<2$ et pour $\handwritten{x}$ et $\handwritten{x}_0$ dans un voisinage assez petit U de $\handwritten{a}$, $\handwritten{y}_0$ dans un voisinage assez petit V de $\handwritten{b}$ et telle que $\handwritten{v}(0;\handwritten{x}_0,\handwritten{x},\handwritten{y}_0)=\handwritten{y}_0$. Si $\handwritten{v}=(v_j)$, on a donc
\[
\frac{dv_j}{dt}=\handwritten{\sum_{i=1}^p}U_{ij}(\handwritten{x}_0+t(\handwritten{x}-\handwritten{x}_0);\handwritten{v}(t))(x_i-x_{i0})\qquad(1\handwritten{\leq}j\handwritten{\leq}q)\tag*{(26)}
\]
On déduit de là que
\[
\begin{aligned}
\frac{d}{dt}\left(\frac{\handwritten{\partial}v_j}{\handwritten{\partial}x_k}\right)&=\frac{\handwritten{\partial}}{\handwritten{\partial}x_k}\left(\frac{dv_j}{dt}\right)=U_{kj}(\handwritten{x}_0+t(\handwritten{x}-\handwritten{x}_0);\handwritten{v}(t))+\\
&\quad+t\handwritten{\sum_i}(x_i-x_{i0})\frac{\handwritten{\partial}U_{ij}}{\handwritten{\partial}x_k}(\handwritten{x}_0+t(\handwritten{x}-\handwritten{x}_0);\handwritten{v}(t))\\
&\quad+\handwritten{\sum_{i,h}}(x_i-x_{i0})\frac{\handwritten{\partial}U_{ij}}{\handwritten{\partial}y_h}(\handwritten{x}_0+t(\handwritten{x}-\handwritten{x}_0);\handwritten{v}(t))\frac{\handwritten{\partial}v_{\unsure{n|h}}}{\handwritten{\partial}x_k}
\qquad(1\handwritten{\leq}j\handwritten{\leq}q,\ 1\handwritten{\leq}k\handwritten{\leq}p)
\end{aligned}
\]
\scanpage{41}
\mspage{27}
% archive sheet number: 40

\[
D\handwritten{u}(x,x_0,\handwritten{y}_0)=\handwritten{f}(x,\handwritten{u}(x,x_0,\handwritten{y}_0))\tag*{(20)}
\]
(D \uline{dériva tion par rapport à} x, $\handwritten{f}=(f_i)$), \uline{et la condition}
\[
\handwritten{u}(x_0,x_0,\handwritten{y}_0)=\handwritten{y}_0.\tag*{(21)}
\]

\uline{En outre, il existe un voisinage} $V_1\handwritten{\subset}V$ \uline{de} $\handwritten{b}$ \uline{tel que, pour tout} $\handwritten{y}_0\handwritten{\in}V_1$ \uline{et tout couple de points} $x,x_0$ \uline{de} U, \uline{l'équation} $\handwritten{y}_0=\handwritten{u}(x_0,x,\handwritten{y})$ \uline{ait dans} V \uline{une seule ra-cine} $\handwritten{y}=\handwritten{u}(x,x_0,\handwritten{y}_0)$.

Posons en effet $x-x_0=t$, $\handwritten{y}-\handwritten{y}_0=\handwritten{z}$ ; l'équation différentielle $\handwritten{y}'=\handwritten{f}(x,\handwritten{y})$ devient $\handwritten{z}'=\handwritten{f}(x_0+t,\handwritten{y}_0+\handwritten{z})$, et le second membre est holomorphe pour $(t,\handwritten{z})$ dans un voisinage de 0 dans $K^{p+1}$ et $(x_0,\handwritten{y}_0)$ dans un voisinage de $(a,\handwritten{b})$ ; la condition (21), pour une intégrale $\handwritten{v}$ de cette équation, s'écrit $\handwritten{v}(0;x_0,\handwritten{y}_0)=0$. On est donc dans les conditions d'application du th.2, d'où la première partie de la proposition. D'autre pa\handwritten{-}rt, si U a été pris assez petit, il existe un voisinage $V_1\handwritten{\subset}V$ de $\handwritten{b}$ tel que pour tout $x_0\handwritten{\in}U$ et tout $\handwritten{y}_0\handwritten{\in}V_1$, l'image de U par l'application $x\handwritten{\to u}(x,x_0,\handwritten{y}_0)$ est contenue dans V car si M est le maximum de $\handwritten{\|f}(x,\handwritten{y})\handwritten{\|}$ dans $U\handwritten{\times}C$, on a en vertu du th. des accroissements finis, $\handwritten{\|u}(x,x_0,\handwritten{y}_0)-\handwritten{y}_0\handwritten{\|}=\handwritten{\|u}(x,x_0,\handwritten{y}_0)-\handwritten{u}(x_0,x_0,\handwritten{y}_0)\handwritten{\|\leq}M|x-x_0|$ quels que soient x et $x_0$ dans U. Si $\handwritten{y}=\handwritten{u}(x,x_0,\handwritten{y}_0)\handwritten{\in}V$, la fonction $s\handwritten{\to u}(s,x,\handwritten{y})$ est donc définie dans U et prend la valeur $\handwritten{y}$ au point x ; elle est donc égale à $\handwritten{u}(s,x_0,\handwritten{y}_0)$ en raison du th.2, et on a donc $\handwritten{u}(x_0,x,\handwritten{y})=\handwritten{u}(x_0,x_0,\handwritten{y}_0)=\handwritten{y}_0$. D'ailleurs, si $\handwritten{w\in}V$ est tel que $\handwritten{y}_0=\handwritten{u}(x_0,x,\handwritten{w})$, l'intégrale $s\handwritten{\to u}(s,x,\handwritten{w})$ qui est définie dans U, prend la valeur $\handwritten{y}_0$ au point $x_0$,\unsure{-.|-} donc est égale à $\handwritten{u}(s,x_0,\handwritten{y}_0)$ ; en particulier, elle prend la valeur $\handwritten{y}$ au point x, ce qui montre que $\handwritten{y}=\handwritten{w}$ et achève la démonstration.
\scanpage{43}
\mspage{24}
% archive sheet number: 42

\[
\left\{\begin{aligned}
m-1&=n+\handwritten{\sum_{j=1}^p\sum_{k=1}^{r_j}}s_{kj}\\
\handwritten{\mu_\ell}&=\handwritten{\unsure{\lambda|\Lambda}_\ell}+\handwritten{\sum_{j=1}^p\sum_{k=1}^{r_j}\unsure{\nu|v}}_{kj}\qquad(1\handwritten{\leq\ell\leq}q)
\end{aligned}\right.\tag*{(16)}
\]
On tire de ces conditions que
\[
m+\handwritten{\sum_{\ell=1}^q\mu_\ell}=1+n+\handwritten{\sum_{\ell=1}^q\lambda_\ell}+\handwritten{\sum_{j=1}^p\sum_{k=1}^{r_j}}(s_{kj}+\handwritten{\sum_{\ell=1}^q\nu}_{kj\handwritten{\ell}})
\]
et par suite, pour tout système d'exposants $(s_{kj},\handwritten{\nu}_{kj1},..,\handwritten{\nu}_{kjq})$ qui figure aux seconds membres de (16), on a $s_{kj}+\handwritten{\sum_{\ell=1}^q\nu}_{kj\handwritten{\ell}}<m+\handwritten{\sum_{\ell=1}^q\mu_\ell}$ ; en d'autres termes, et comme K est de caractéristique 0, les relations (15) déterminent les coefficients des termes de degré \uline{total} t dans $u_i$ en fonction de ceux des termes de degré total $<t$ dans les $u_j$ $(1\handwritten{\leq}j\handwritten{\leq}p)$ ; par récurrence sur t, on voit donc que les équations (13) déterminent les coefficients $h_i(m;\handwritten{\mu}_1,..,\handwritten{\mu}_q)$ de fa\handwritten{-}çon \uline{unique}, ce qui montre déjà que si le problème est possible, il n'a qu'une seule solution. En outre, si les $h_i(m;\handwritten{\mu}_1,..,\handwritten{\mu}_q)$ sont déterminés par récurrence à partir des relations (15), les séries entières ayant ces coefficients seront solutions du problème \uline{pourvu qu'elles soient sommables dans un\unsure{voisina-ge|voisinage} de l'origine} dans $K^{q+1}$ ; . En effet, s'il en est ainsi, ces séries sont fonctions continues dans un voisina ge de l'origine, donc on peut trouver un tel voisinage $U\handwritten{\times}W$ dont l'image par $(x;z_1,..,z_q)\handwritten{\to}(u_i(x;z_1,\ldots,z_q))$ soit contenue dans P, et tel que les séries $u_i$ soient substituables dans (14) pour $x\handwritten{\in}U$, $(z_1,\ldots,z_q)\handwritten{\in}W$ (\handwritten{\S}1,th.2) ; l'expression de la dérivée d'une série entière et les relations (15) montreront alors que les $u_i$ sont bien solutions du système (12) dans $U\handwritten{\in}W$.

Pour prouver la \unsure{somm bilité|sommabilité} des séries $u_i$, considérons pour chacune des séries (14), une série majorante, dont nous désignerons les coefficients par $C_i(n;r_1,..,r_p;\handwritten{\lambda}_1,..,\handwritten{\lambda}_q)$ ; déterminons alors les coefficients $\unsure{\mathbb{H}|H}_i(m;\handwritten{\mu}_1,\ldots,\handwritten{\mu}_q)$ par les conditions
\scanpage{45}
\mspage{23}
% archive sheet number: 44

\[
Du_i(x;z_1,..,z_q)=f_i(x;u_1(x;z_1,..,z_q),..,u_p(x;z_1,..,z_q);z_1,..,z_q)\tag*{(12)}
\]
(D \uline{dérivation par rapport à} x, $1\handwritten{\leq}i\handwritten{\leq}p$) \uline{et}
\[
u_i(a;z_1,..,z_q)=b_i\qquad(1\handwritten{\leq}i\handwritten{\leq}p)\quad\text{\uline{pour tout}}\quad(z_1,..,z_q)\handwritten{\in}W\tag*{(13)}
\]

Par translation, on peut se ramener au cas où $a=b_j=c_k=0$ pour $1\handwritten{\leq}j\handwritten{\leq}p$ et $1\handwritten{\leq}k\handwritten{\leq}q$ ; chacune des fonctions $f_i$ est alors égale à la somme d'une série entière sommable dans un polycylindre ouvert $B\handwritten{\times}P\handwritten{\times}Q$, où $B\handwritten{\subset}K$, $P\handwritten{\subset}K^p$ et $Q\handwritten{\subset}K^q$
\[
\begin{aligned}
f_i(x;y_1,..,y_p;z_1,..,z_q)&=\handwritten{\sum}\unsure{g|e|c}_i(n;r_1,..,r_p;\handwritten{\unsure{\lambda|\Lambda}}_1,..,\handwritten{\lambda}_q)\mathbin{.}\\
&\quad\mathbin{.}x^n y_1^{r_1}..y_p^{r_p}z_1^{\handwritten{\lambda}_1}z_2^{\handwritten{\lambda}_2}\ldots z_q^{\handwritten{\lambda}_q}.
\end{aligned}\tag*{(14)}
\]

Pour établir l'existence des fonctions $u_i$, nous devons prouver qu'il existe p séries entières sans terme indépendant de x
\[
u_i(x;z_1,..,z_q)=\handwritten{\sum}h_i(m;\handwritten{\mu}_1,..,\handwritten{\mu}_q)x^m z_1^{\handwritten{\mu}_1}\ldots z_q^{\handwritten{\mu}_q}
\]
sommables dans un polydylindre $U\handwritten{\times}W$ de $K^{q+1}$, substituables dans les séries entières (14) et satisfaisant à (12). Si le problème est possible, les séries entières par rapport à x et aux $z_{\handwritten{\ell}}$ $(1\handwritten{\leq\ell\leq}q)$ obtenues par cette substitution dans les deux membres de (12) doivent avoir mêmes coefficients, ce qui donne les conditions
\[
\left.\begin{aligned}
h_i(0;\handwritten{\mu}_1,..,\handwritten{\mu}_q)&=0\\
mh_i(m;\handwritten{\mu}_1,\ldots,\handwritten{\mu}_q)&=\handwritten{\sum_{n,(r_j),(\lambda_\ell)}}g_i(n;r_1,..,r_p;\handwritten{\lambda}_1,\ldots,\handwritten{\lambda}_q)\mathbin{.}\\
&\quad\mathbin{.}\handwritten{\prod_{j=1}^p\prod_{k=1}^{r_j}}h_i(s_{kj};\handwritten{\unsure{\nu|v}}_{kj1},\ldots,\handwritten{\nu}_{kjq})
\end{aligned}\right\}\tag*{(15)}
\]
pour tout entier $m\handwritten{>}0$ et tout système d'entiers $\handwritten{\mu_\ell\geq}0$ $(1\handwritten{\leq\ell\leq}q)$ ; la sommation du second membre des relations (15) doit être étendue à tous les systèmes d'entiers $\handwritten{\geq}0$, $n,r_j,\handwritten{\lambda_\ell}$, $s_{kj}$ et $\handwritten{\nu}_{kj\handwritten{\ell}}$ satisfaisant aux conditions
\scanpage{47}
\mspage{22}
% archive sheet number: 46

COROLLAIRE.- \uline{Soit} $\handwritten{u}(z_1,..,z_p)$ \uline{une fonction holomorphe dans une partie ouverte} $A$ \uline{de} $K^p$, \uline{à valeurs dans} $K^p$. \uline{Pour tout point} $\handwritten{a}=(a_1,..,a_p)$ \uline{de} $A$ \uline{tel que la différentielle} $\handwritten{\unsure{d.u(a)|du(a)}}$ \uline{soit inversible il existe un voisinage ouvert} $U$ \uline{de} $\handwritten{a}$ \uline{tel que la restriction de} $\handwritten{u}$ \uline{à} $U$ \uline{soit un homéomorphisme de} $U$ \uline{sur un voisinage ouvert} $V$ \uline{de} $\handwritten{u(a)}$, \uline{et que l'homéomorphisme réciproque soit holomorphe dans} $V$.

\margin{\handwritten{\figure{\begin{tikzpicture}[x=0.10cm,y=0.10cm,baseline=0pt]
\fill (0,0) -- (-0.15,1.8) -- (-1.15,1.5) -- (-0.8,0.4) -- cycle;
\draw[line width=0.55pt] (-0.4,1.35) -- (-1.6,1.1) -- (-2.9,0.9);
\draw[line width=0.45pt] (-0.55,0.6) .. controls (-1.3,1.05) and (-1.8,0.2) .. (-2.15,0.55);
\end{tikzpicture}}}}
Lorsque $K=\handwritten{\unsure{\mathbb{R}|R}}$ et que $A$ est un intervalle ouvert dans $\handwritten{R}$, le fait que $u'(x)\handwritten{\ne}0$ dans $A$ entraîne que $u$ est un homéomorphisme de $A$ \uline{tout entier} sur $u(A)$ puisque $u$ est monotone dans $A$ (\uline{Fonct.var.réelle}, chap.I,\S\,2 et \uline{Top.gén.} chap.IV,\S\,2). Mais lorsque $K=\handwritten{\mathbb{C}}$, $u'(x)$ peut être $\handwritten{\ne}0$ en tout point d'un domaine $A\handwritten{\subset C}$ sans que $u$ soit un homéomorphisme de $A$ sur $u(A)$. Par exemple, si $u(x)=x^n$, on a $u'(x)\handwritten{\ne}0$ pour $x\handwritten{\ne}0$, mais si $A$ est le disque $\handwritten{|x-1|<\frac{1}{2}}$, $u$ n'est pas un homéomorphisme de $A$ sur $u(A)$ dès que $n$ est assez grand, car il existe des points $x\handwritten{\ne}1$ dans $A$ tels que $x^n=1$.

Il faudrait faire ici également le théorème de \uline{dépendance} des fonctions holomorphes à jacobien identiquement nul.

\uline{3}. \uline{Intégration des équations différentielles}.

THEOREME 2.- \uline{Soit} $K$ \uline{un corps valué complet de caractéristique} $0$. \uline{Soient} \margin{\handwritten{\unsure{-}}}$f_i(x;y_1,..,y_p;z_1,..,z_q)$ $(1\handwritten{\leq}i\handwritten{\leq}p)$ $p$ \uline{fonctions holomorphes scalaires définies dans une partie ouverte} $A$ \uline{de} $K^{p+q+1}$. \uline{Pour tout point} $(a;b_1,..,b_p;c_1,..,c_q)$ \uline{de} $A$, \uline{il existe un voisinage} $U$ \uline{de} $a$ \uline{dans} $K$, \uline{un voisinage} $V$ \uline{de} $(b_1,..,b_p)$ \uline{et un voisinage} $W$ \uline{de} $(c_1,..,c_q)$ \uline{tels que} $U\handwritten{\times}V\handwritten{\times}W\handwritten{\subset}A$, \uline{et un système et un seul de} $p$ \uline{fonctions scalaires} $u_i(x;z_1,..,z_q)$, \uline{holomorphes dans} $U\handwritten{\times}W$ \uline{telles que l'image de} $U\handwritten{\times}W$ \uline{par l'application} $(x;z_1,..,z_q)\handwritten{\to}(u_i(x;z_1,..,z_q))$ \uline{soit contenue dans} $V$ \uline{et qui satisfont dans} $U\handwritten{\times}W$ \uline{aux conditions}
